Twelve designs outside the training set were specified at the end of Stage 6. Their strength, modulus, failure strain, work to failure, damage localisation and fracture-mode class were predicted and written to a timestamped, hashed file before any of them was simulated (file written 2026-09-05 21:20:34, SHA-256 in Table~\ref{tab:holdout}; earliest holdout record 2026-09-05 21:52:25). The predictor is a documented, deliberately simple data-driven rule: a nearest-neighbour regression in descriptor space (porosity, minimum load-bearing solid fraction, ligament width, anisotropy, hierarchy level, disorder) over the Stage 1--6 records, with the net-section strength assigned by ligament class (Sec.~\ref{sec:mech}). Its failures therefore point at missing physics rather than at fitting noise. Table~\ref{tab:holdout} and Fig.~\ref{fig:holdout} compare predictions with observations.

Where the net-section rule applies unchanged, the predictions are accurate: the 90-degree strut lattice at 30\% porosity (14.4 predicted, 14.4 observed), the 50\,\AA{} crack in a 150\,\AA{} cell (16.3 vs 16.1\,N/m, confirming the lattice-trapping plateau at a crack length and cell size not used before), the three-level mesh with wide veins (11.7 vs 12.3), the jittered mesh with strength dispersion (12.3 vs 11.9), the 24\,\AA{} mesh at 30\% porosity (10.4 vs 9.9) and the staggered square mesh (19.3 vs 17.7\,N/m) all fall within 10\%. Over the twelve designs the median absolute error of the strength prediction is 14\% (Stage 4 discriminating experiments: 6.5\%), of the failure strain 17\%, of the work to failure 20\%, of the damage localisation 10\% and of the 2D modulus 11\%.

The five larger misses are the scientifically informative part of the test:
\begin{itemize}
\item \emph{Slit array rotated by 20 degrees} (predicted 23.9, observed 8.9\,N/m). The predictor treated a nearly load-parallel slit array like a parallel one. In the simulation the slit tips of neighbouring rows overlap along the load direction, and the ligaments between them fail by en-echelon crack coalescence (Fig.~\ref{fig:comparison_slit_orientation_echelon}), a mechanism that did not occur anywhere in the training set. The strength lies far below both the parallel (25.4\,N/m) and the 45-degree (18.5\,N/m) arrays of Stage 2 and closer to the perpendicular array (4.0\,N/m), although the sheet then fails progressively along the echelon paths and retains about 8\,N/m to 19\% strain. Load-path alignment is therefore not a monotonic function of angle; it is lost as soon as the tips of adjacent slit rows overlap along the load.
\item \emph{Load-parallel veins combined with load-elongated fine pores} (predicted 13.3, observed 19.4\,N/m; Fig.~\ref{fig:comparison_aligned_hierarchy_synergy}). This is the strongest nested design of the campaign; it failed in a few steps (19.4 to 11.8\,N/m at the first event) and its work to failure (2.0\,J/m$^2$) equals that of the best Stage 5 nested design, which had square veins. The two orientation effects act on different length scales and their benefits add; the nearest-neighbour predictor averaged over nested designs with and without alignment and could not represent the combination.
\item \emph{Parallel slits in armchair-oriented graphene} (predicted 24.0, observed 20.2\,N/m). The predictor had no lattice-orientation input. The net-section rule with the armchair intrinsic strength of Stage 1 (30.7 instead of 38.8\,N/m) predicts $25.4 \times 30.7/38.8 = 20.1$\,N/m, within 1\% of the observation: the ligament strength inherits the chirality dependence of the pristine sheet.
\item \emph{Flaw halo of three rings around a 30\,\AA{} hole} (predicted 13.1, observed 8.9\,N/m) and \emph{graded mesh with a weak band} (15.4 vs 10.6\,N/m). Both were predicted from cell-averaged descriptors; both are controlled by their weakest cross-section. The rings turn the hole into a wider effective flaw bounded by thin, curved ligaments, and the graded sheet fails in its most porous band. Strength is a minimum, not an average, over cross-sections, as the Stage 5 gradient series had already indicated.
\item \emph{Voronoi network with 35 cells and regularity 0.8} (predicted 12.1, observed 7.3\,N/m). A new seed at a cell count outside the training range produced the weakest holdout. The miss is consistent with the seed-to-seed scatter of random networks measured in Stage 6 (9.2 to 14.4\,N/m for three seeds of the 25-cell network at regularity 0.6, Fig.~\ref{fig:seeds}) and with the disorder penalty of hypothesis H-D, not with a new mechanism.
\end{itemize}

The fracture-mode class was predicted correctly in 5 of 12 cases in the strict three-class sense (single avalanche; stepwise with few events; progressive with load retention) and in 10 of 12 cases in the binary sense (single avalanche versus multi-event). All five strict misses are confusions between the two multi-event classes. The two binary misses are the 50\,\AA{} crack, which advanced bond by bond before the instability so that lattice trapping made it stepwise rather than abrupt, and the armchair slit array, whose wide ligaments failed one after another (progressive rather than abrupt). Compared with Stage 4 (median strength error 6.5\%, mode accuracy 80\%) the holdouts were harder by design: they combined variables, extended ranges, and included one geometry, the rotated slit array, chosen to probe the limit of the alignment rule.

The holdout test changes the conclusions in four ways: (i) the net-section rule holds as a minimum over cross-sections with a chirality-dependent ligament strength; (ii) alignment benefits at different scales add, which makes the aligned two-level mesh the best nested architecture found; (iii) the alignment rule has a sharp angular limit set by tip overlap (en-echelon coalescence); (iv) data-driven interpolation over descriptors, even with the campaign's mechanistic classes built in, fails exactly where a mechanism changes, which is the argument for keeping the validated simulation in the loop rather than replacing it by a surrogate.
